% 第 4 章 §4.1–§4.2.1 前段（印刷页 165–171）
\chapter{Riemann 面上的 Higgs 丛与局部系统}
\label{chap:04}

\noindent Richard A. Wentworth\footnote{R.W. 部分受 NSF 基金 DMS--1037094 与 DMS--1406513 资助。作者还感谢 NSF 基金 DMS 1107452、1107263、1107367「RNMS: 几何结构与表示簇」（GEAR 网络）的支持。}

\section{前言}

这些讲义基于 2012 年 3 月 26--30 日在巴塞罗那数学研究中心（Centre de Recerca Matem\`atica）举行的「几何与物理第三届国际学校」上的报告。该学校四个系列讲座的目标是快速介绍 Higgs 丛、表示簇与数学物理。虽然这些主题的范围极广，但本讲义的范围要小得多。这里涵盖的主要论题是：

\begin{enumerate}[label=(\roman*), leftmargin=2.6em, itemsep=0.2em]
\item Riemann 面上 Higgs 丛的 Hitchin--Kobayashi--Simpson 对应；
\item 联系 Higgs 丛模空间与基本群半单表示的 Corlette--Donaldson 定理；
\item oper 模空间的描述，及其与全纯微分方程组、Higgs 丛以及 Eichler--Shimura 同构的关系。
\end{enumerate}

这些主题在文献中已被广泛处理。我试图把关键思想压缩成一份所需背景尽可能少的陈述。关于第一项，我给出 Hitchin--Simpson 定理（定理 4.2.17）的完整证明，其中综合了自 Hitchin 的开创性论文 [35] 以来出现的各种技巧。在 Riemann 面的情形，可以仿照 Donaldson 在 [18] 中对 Narasimhan--Seshadri 定理的证明，给出任意秩的直接证明（这样的证明曾在 [56] 中被提议）。此外，Yang--Mills--Higgs 流可用来抽取具有所需性质的最小化序列。类似的想法用在与 Corlette--Donaldson 证明等变调和映射存在性（定理 4.3.14）时。事实上，我在本讲义中一直试图展示这两个基本结果之证明的平行结构。两个流的连续性，是半稳定 Higgs 丛的模空间的等变上同调与表示的模空间之间关系的关键。乍看之下，上面列表中的最后一项与其余各项相当不同，但它以或许尚未被完全理解的方式与之深刻相关。oper [3] 在几何 Langlands 纲领 [24] 的文献中扮演重要角色。我的意图是给出关于 oper 及其与微分方程、Higgs 丛之关系的基本事实的相当完整的证明（另见 [62]）。

由于时间和篇幅的限制，我不得不略去这一主题的许多重要方面。有两处特别值得一提。第一，我只处理向量丛，不考虑具有更一般结构群的主丛。例如，没有讨论到复李群的各种实形式的表示。由于这所入门学校的其他一些讲座将详细处理这一主题，我希望这一省略不会太严重。第二，我只处理闭 Riemann 面，不考虑标记点处的额外「抛物」结构。在某种意义下，这忽略了经典全纯微分方程文献核心的一个重要方面（参看 [7, 58]）。然而，为了介绍模空间的整体结构，我觉得先处理闭曲面的情形更好。尽管该领域目前的许多研究正朝着上述两种推广进行，这些主题留作进一步阅读。

我试图在本讲义中为关键结果给出参考文献。任何遗漏或不正确的归属，都完全出于我对这一极其丰富浩瀚的文献的无知，为此我致以诚挚的歉意。此外，这里给出的证明没有任何原创性主张。强烈建议任何想学习 Higgs 丛的人仔细阅读 Carlos Simpson 在这一主题上的奠基性贡献（见 [56, 57, 59, 60, 61]）。此外，Corlette [11]、Donaldson [18, 20]，以及当然还有 Hitchin [35, 36, 38] 的原始文章都是不可或缺的。最后，我还要提到较新的综述文章 [9, 10, 27]，它们特别处理到一般李群的表示的情形。

我感谢组织者 Luis \'Alvarez-C\'onsul、Peter Gothen 与 Ignasi Mundet i Riera 邀请我做这些报告，也感谢 CRM 的热情接待。还要感谢 Bill Goldman、Fran\c{c}ois Labourie、Andy Sanders 与 Graeme Wilkin 就这里介绍的主题所作的讨论，感谢 Beno\^{\i}t Cadorel 找出若干笔误。匿名审稿人也提出了非常有用的建议，我对此深表感激。

\subsection*{记号}

\begin{itemize}[leftmargin=1.6em, itemsep=0.15em]
\item $X$ = 亏格 $g\ge2$ 的紧 Riemann 面；
\item $\pi=\pi_1(X,p)$ = $X$ 的基本群；
\item $\mathbb{H}$ = $\C$ 中的上半平面；
\item $\mathcal{O}=\mathcal{O}_X$ = $X$ 上全纯函数芽层；
\item $K=K_X$ = $X$ 的典范层；
\item $E$ = $X$ 上的复向量丛；
\item $H$ = $E$ 上的埃尔米特度量；
\item $\nabla$ = $E$ 上的联络；
\item $A$（或 $d_A$）= $(E,H)$ 上的酉联络；
\item $\mathcal{C}_E$ = 秩 $n$ 丛 $E$ 上联络的空间；
\item $\mathcal{A}_E$ = $E$ 上酉联络的空间；
\item $\mathcal{B}_E$ = Higgs 丛的空间；
\item $\mathcal{B}^{ss}_E$ = 半稳定 Higgs 丛的空间；
\item $\mathcal{G}_E$（相应地 $\mathcal{G}_E^\C$）= 酉（相应地复）规范群；
\item $\bar\partial_E$ = $E$ 上的 Dolbeault 算子，等价于一个全纯结构；
\item $(\bar\partial_E,H)$ = Chern 连接；
\item $\mathcal{E}$ = 全纯丛 $(E,\bar\partial_E)$ 的全纯截面芽层；
\item $\mathfrak{g}_E$ = $E$ 的反埃尔米特自同态丛；
\item $\operatorname{End}E=\mathfrak{g}^\C_E$ = $E$ 的自同态丛；
\item $V$ = $X$ 上的局部系统；
\item $V_\rho$ = 联系于表示 $\rho\colon\pi\to GL_n(\C)$ 的局部系统；
\item $\underline{R}$ = 以环 $R$ 为模型的局部常值层；
\item $L^p_k$ = $k$ 阶导数属于 $L^p$ 的函数/截面构成的 Sobolev 空间；
\item $C^{k,\alpha}$ = $k$ 阶导数 H\"older 连续（指数 $\alpha$）的函数/截面空间。
\end{itemize}

\section{Dolbeault 模空间}

\subsection{Higgs 丛}

\paragraph{全纯丛与稳定性} 本讲义中 $X$ 总表示亏格 $g\ge2$ 的闭 Riemann 面，$E\to X$ 是复向量丛。我们先讨论复向量丛的基本微分几何。这方面好的参考文献是 Kobayashi 的书 [44] 与 Griffiths--Harris [26]。$E$ 上的全纯结构等价于一个 $\bar\partial$-算子的选取，即一个 $\C$-线性映射
\[
\bar\partial_E\colon\Omega^0(X,E)\longrightarrow\Omega^{0,1}(X,E),
\]
满足 Leibniz 法则：对函数 $f$ 与 $E$ 的截面 $s$，有 $\bar\partial_E(fs)=\bar\partial f\otimes s+f\bar\partial_E s$。事实上，若 $\{s_i\}$ 是全纯丛的局部全纯标架，则 Leibniz 法则唯一确定底复向量丛上的 $\bar\partial$-算子。反之，由于 Riemann 面上没有可积性条件，给定如上定义的 $\bar\partial$-算子，总能找到局部全纯标架（参看 [2, \S5]）。当要指明全纯结构时，我们写 $(E,\bar\partial_E)$。我们还用 $\mathcal{E}$ 记 $(E,\bar\partial_E)$ 的全纯截面芽层。有时我们混用术语，把 $\mathcal{E}$ 称为全纯丛。

若 $S\subset E$ 是全纯子丛，商丛为 $Q$，则光滑分裂 $E=S\oplus Q$ 使我们能把 $\bar\partial$-算子写成
\begin{equation}
\bar\partial_E=\begin{pmatrix}\bar\partial_S & \beta\\ 0 & \bar\partial_Q\end{pmatrix},
\label{eq:4.1}\tag{4.1}
\end{equation}
其中 $\beta\in\Omega^{0,1}(X,\operatorname{Hom}(Q,S))$ 称为第二基本形式。$E$ 上的埃尔米特度量 $H$ 给出正交分裂。此时子丛 $S$ 由其正交投影算子 $\pi$ 决定，$\pi$ 是 $E$ 的自同态，满足：(i) $\pi^2=\pi$；(ii) $\pi^*=\pi$；(iii) $\operatorname{tr}\pi$ 为常数。$S\subset E$ 全纯这一陈述等价于附加条件 (iv) $(I-\pi)\bar\partial_E\pi=0$。注意 (i) 与 (iv) 蕴涵 (iii)，且 $\beta=-\bar\partial_E\pi$。因此，$E$ 的全纯子丛与埃尔米特丛 $E$ 上满足 (i)、(ii)、(iv) 的自同态 $\pi$ 之间有一一对应。我应指出，把这一关于全纯子层的描述推广到高维，是 Uhlenbeck--Yau [65] 的关键思想。

$E$ 上的联络 $\nabla$ 是一个 $\C$-线性映射 $\nabla\colon\Omega^0(X,E)\to\Omega^1(X,E)$，满足 Leibniz 法则：$\nabla(fs)=df\otimes s+f\nabla s$。给定埃尔米特度量 $H$，若联络保持 $H$，即
\begin{equation}
d\langle s_1,s_2\rangle_H=\langle d_As_1,s_2\rangle_H+\langle s_1,d_As_2\rangle_H,
\label{eq:4.2}\tag{4.2}
\end{equation}
则称之为\emph{酉}联络（此时我们总把它记作 $A$ 或 $d_A$）。

联络 $\nabla$ 的曲率是 $F_\nabla=\nabla^2$（更精确地写作 $\nabla\wedge\nabla$）。若 $\mathfrak{g}_E$ 表示 $E$ 的反埃尔米特自同态丛，$\mathfrak{g}^\C_E$ 是它的复化，则对酉联络有 $F_A\in\Omega^2(X,\mathfrak{g}_E)$，一般情形 $F_\nabla\in\Omega^2(X,\mathfrak{g}^\C_E)$。

\begin{remark}
我们主要处理那些在行列式丛上诱导固定联络的丛上的联络。例如，它们对应于到 $SL_n$ 而非 $GL_n$ 的表示。此时，丛 $\mathfrak{g}_E$ 与 $\mathfrak{g}^\C_E$ 应取为无迹自同态的丛。
\end{remark}

最后注意，联络总可通过取其 $(0,1)$ 部分而诱导一个 $\bar\partial$-算子。反之，一个 $\bar\partial$-算子给出唯一的酉联络，称为 Chern 连接，我们有时记作 $d_A=(\bar\partial_E,H)$。$X$ 上的复结构把 $\Omega^1(X)$ 分裂为 $(1,0)$ 与 $(0,1)$ 部分，从而也分裂联络。我们分别记之为 $d'_A$ 与 $d''_A$。于是对 $d_A=(\bar\partial_E,H)$，有 $d''_A=\bar\partial_E$，而 $d'_A$ 由 $\partial\langle s_1,s_2\rangle_H=\langle d'_As_1,s_2\rangle_H$ 对任意全纯截面 $s_1,s_2$ 确定。往后，若无混淆之虞，我大多会从记号中省略 $H$。

\begin{example}
设 $L$ 是有埃尔米特度量 $H$ 的全纯线丛。对局部全纯标架 $s$，写 $H_s=|s|^2$。则 $F_{(\bar\partial_L,H)}=\bar\partial\partial\log H_s$，且右端与标架的选取无关。
\end{example}

全纯线丛在 $X$ 的覆盖的各开集上的一族局部平凡化，其转移函数给出取值于处处非零全纯函数芽层 $\mathcal{O}^*$ 的 $1$-上循环。线丛的同构类之集于是是 $H^1(X,\mathcal{O}^*)$。回忆紧 Riemann 面上每个全纯线丛都有亚纯截面。这给出全纯线丛（在张量积下）的范畴与除子 $D=\sum_{x\in X}m_xx$ 的线性等价类（连同其加法结构）之间的等价（这里 $m_x\in\Z$ 除有限多个 $x$ 外为零）。我们用 $\mathcal{O}(D)$ 记如此联系于 $D$ 的线丛。此外，除子有次数 $\deg D=\sum_{x\in X}m_x$。我们把它定义为 $\mathcal{O}(D)$ 的次数。

另一方面，由指数正合列
\[
0\longrightarrow\Z\longrightarrow\mathcal{O}\xrightarrow[\quad f\mapsto e^{2\pi if}\quad]{}\mathcal{O}^*\longrightarrow0,
\]
得上同调的长正合列：
\[
0\longrightarrow H^1(X,\Z)\longrightarrow H^1(X,\mathcal{O})\longrightarrow H^1(X,\mathcal{O}^*)\xrightarrow{c_1}H^2(X,\Z)\longrightarrow0.
\]
$X$ 的基本类把 $H^2(X,\Z)$ 等同于 $\Z$，而标准的练习是证明在此等同下 $\deg(D)=c_1(\mathcal{O}(D))$。对全纯向量丛 $E$，我们定义 $\deg E:=\deg\det E$。注意次数是拓扑的，即不依赖全纯结构，只依赖底复丛 $E$。由 Chern--Weil 理论，对 $E$ 上任一埃尔米特度量 $H$ 有
\begin{equation}
c_1(E)=\left\langle\frac{\sqrt{-1}}{2\pi}\operatorname{tr}F_{(\bar\partial_E,H)}\right\rangle=\left\langle\frac{\sqrt{-1}}{2\pi}F_{(\bar\partial_{\det E},\det H)}\right\rangle .
\label{eq:4.3}\tag{4.3}
\end{equation}

Riemann 面上的复向量丛按秩与次数作拓扑分类。我们还会用到丛的\emph{斜率}（或归一化次数），它定义为比值 $\mu(E)=\deg E/\operatorname{rank}E$。

若线丛 $L=\mathcal{O}(D)$ 有非零全纯截面，则由 $D$ 线性等价于有效除子（即对一切 $x$ 有 $m_x\ge0$），得 $\deg L\ge0$。由此推出：若 $E$ 是全纯向量丛，$S\subset E$ 是子层且 $\operatorname{rank}S=\operatorname{rank}E$，则 $\deg S\le\deg E$。事实上，该假设蕴涵 $\det E\otimes(\det S)^*$ 有非零全纯截面。我们以后会用到这一事实。注意在上述情形 $Q=E/S$ 是挠层。一般地，对全纯向量丛的任一子层 $S\subset E$，$S$ 含于 $E$ 的唯一确定的全纯子丛 $S'$ 中，后者称为 $S$ 的\emph{饱和化}。它由诱导映射 $E\to Q/\operatorname{Tor}(Q)\to0$ 的核得到。由这一讨论我们断定 $\deg S$ 不超过其饱和化的次数 $\deg S'$。

设 $\omega$ 是 $X$ 上某个共形度量对应的凯勒形式。它在全文中固定，为方便起见我们归一化使
\[
\int_X\omega=2\pi .
\]
缩并 $\Lambda\colon\Omega^2(X)\to\Omega^0(X)$ 由对任意函数 $f$ 令 $\Lambda(f\omega)=f$ 定义。对埃尔米特全纯丛 $E$ 的全纯子丛 $S$（投影算子为 $\pi$），有如下有用的公式，它容易由 (4.3) 直接计算得到：
\begin{equation}
\deg S=\frac{1}{2\pi}\int_X\operatorname{tr}\left(\pi\sqrt{-1}\Lambda F_{(\bar\partial_E,H)}\right)\omega-\frac{1}{2\pi}\int_X|\beta|^2\omega .
\label{eq:4.4}\tag{4.4}
\end{equation}

\begin{definition}
若对一切全纯子丛 $S\subset E$（$0<\operatorname{rank}S<\operatorname{rank}E$）有 $\mu(S)<\mu(E)$（相应地 $\mu(S)\le\mu(E)$），则称 $E$ 为\emph{稳定}（相应地\emph{半稳定}）的。若 $E$ 是同斜率稳定丛的直和，则称其为\emph{多稳定}的。
\end{definition}

\begin{remark}
线丛平凡地稳定。若 $E$ 是（半）稳定的而 $L$ 是线丛，则 $E\otimes L$ 也是（半）稳定的。
\end{remark}

在给例子之前，回忆\emph{扩张}的概念
\begin{equation}
0\longrightarrow S\longrightarrow E\longrightarrow Q\longrightarrow0 .
\label{eq:4.5}\tag{4.5}
\end{equation}
扩张类是与 (4.5) 联系的长正合列的上边缘映射下恒等自同态的像，
\[
H^0(X,Q\otimes Q^*)\longrightarrow H^1(X,S\otimes Q^*).
\]
注意丛 $E$ 的同构类在伸缩下不变，故扩张类（若非零）应视为射影空间 $\PP(H^1(X,S\otimes Q^*))$ 的元素。于是容易验证：就第二基本形式 $\beta$ 而言，扩张类（在射影意义下）与相应的 Dolbeault 上同调类 $[\beta]\in H^{0,1}_{\bar\partial}(X,S\otimes Q^*)$ 一致。若扩张类为零，则称 (4.5) 分裂。显然这发生当且仅当存在提升投影的单射 $Q\to E$。

\begin{example}
设 $g\ge1$。考虑
\[
0\longrightarrow\mathcal{O}\longrightarrow E\longrightarrow\mathcal{O}(p)\longrightarrow0
\]
型的扩张。它们由 $H^1(X,\mathcal{O}(-p))\cong H^0(X,K(p))^*\cong H^0(X,K)^*$ 参数化，其维数为 $g$。这类任一非分裂扩张都稳定。事实上，若 $L\hookrightarrow E$ 是破坏稳定的线子丛，则 $\deg L\ge1$。诱导映射 $L\to\mathcal{O}(p)$ 不能为零，因为那样由包含 $L\hookrightarrow E$ 它会提升为非零映射 $L\to\mathcal{O}$，而这是不可能的。因此 $L\to\mathcal{O}(p)$ 必为同构。这样的 $L$ 就会分裂该扩张。
\end{example}

若联络的曲率为零，则称之为\emph{平坦}的。若 $\sqrt{-1}\Lambda F_\nabla=\mu$，其中 $\mu$ 是常数（单位矩阵的倍数），则称 $\nabla$ 为\emph{射影平坦}的。注意由我们对面积的归一化，$\mu=\mu(E)$。在 §4.4 中我们将证明 Weil 判别法，即全纯丛 $E$ 何时容许平坦联络（即 $\nabla''=\bar\partial_E$，$F_\nabla=0$）。要求联络是酉的会施加更强的条件。这就是著名的 Narasimhan--Seshadri 结果。

\begin{theorem}[Narasimhan--Seshadri, [49]]\label{thm:4.2.6}
全纯丛 $E\to X$ 容许射影平坦酉联络，当且仅当 $E$ 是多稳定的。
\end{theorem}

在 §4.2.3 中我们将把定理 4.2.6 作为 Higgs 丛的更一般结果（见定理 4.2.17）的特例来证明。

\paragraph{Higgs 场} \emph{Higgs 丛}是一对 $(E,\Phi)$，其中 $E$ 是全纯丛，$\Phi$ 是 $K\otimes\operatorname{End}E$ 的全纯截面。我们有时把 $\Phi$ 视为满足 $\bar\partial_E\Phi=0$ 的 $\Omega^{1,0}(X,\mathfrak{g}^\C_E)$ 的截面。

\begin{definition}
若对一切 $\Phi$-不变的全纯子丛 $S\subset E$（$0<\operatorname{rank}S<\operatorname{rank}E$）有 $\mu(S)<\mu(E)$（相应地 $\mu(S)\le\mu(E)$），则称 $(E,\Phi)$ 为\emph{稳定}（相应地\emph{半稳定}）的。若它是同斜率 Higgs 丛的直和，则称其为\emph{多稳定}的。
\end{definition}

下面是由定义以及斜率在正合列上的可加性质得到的一个简单但有用的推论。
